% Charge d'un condensateur. Résolution analytique
\begin{minipage}{0.55\linewidth}
    Un condensateur initialement déchargé est inséré dans le montage suivante.


    À la date $t=0$, on ferme l'interrupteur $\mathrm{K}$.

    \begin{enumerate}[series=x]
        \item Représenter par une flèche le sens de circulation du courant
              d'intensité $i$ dans le circuit. Identifier l'armature positive du
              condensateur. Représenter par des flèches les tensions $u_{C}$ aux
              bornes du condensateur et $u_{R}$ aux bornes du conducteur
              ohmique, en respectant la convention récepteur.
        \item En appliquant la loi d'additivité des tensions, établir une
              relation~(1) entre $E$, $u_{R}$ et $u_{C}$.
    \end{enumerate}
\end{minipage}
\hfill
\begin{minipage}{0.44\linewidth}
    \begin{figure}[H]
        \centering
        \begin{circuitikz}
            \newdimen\d \d=3cm
            \draw (0,0) coordinate (O)
                to[vsource,v^>=$E$] ++(0,\d)
                to[cute open switch,l=$\mathrm{K}$] ++(\d,0)
                to[R,l=$R$,-*] ++(\d,0) node[right] {$B$}
                to[C,l_=$C$,-*] ++(0,-\d) node[right] {$A$}
                to[short] (O);
        \end{circuitikz}
    \end{figure}
\end{minipage}


\begin{enumerate}[resume=x]
    \item Exprimer $u_{R}$ en fonction de $i$ et l'intensité $i$ en fonction de
          la charge $q$ du condensateur.
    \item En déduire l'expression de $i$ en fonction de la capacité $C$ et de la
          tension $u_{C}$.
    \item À l'aide de la relation~(1), établir l'équation différentielle à
          laquelle obéit $u_{C}$.
    \item Une solution de cette équation différentielle est de la forme~:
          $u_{C}=c+a\times\mathrm{e}^{bt}$.
          \begin{enumerate}
              \item En reportant cette expression dans la relation~(1),
                    déterminer la valeur des constantes $b$ et $c$.
              \item À $t=0$, que vaut la tension $u_{C}$~? En déduire la valeur
                    du coefficient $a$.
              \item En déduire l'expression de $u_{C}$ en fonction du temps.
          \end{enumerate}
\end{enumerate}

\begin{donnees}
    \begin{itemize}
      \item $R=\qty{10}{\kilo\ohm}$~; $C=\qty{2200}{\uF}$~; $E=\qty{6.0}{\V}$.
    \end{itemize}
\end{donnees}


